% Part: applied-modal-logics
% Chapter: epistemic-logic
% Section: relational-models

\documentclass[../../../include/open-logic-section]{subfiles}

\begin{document}

\olfileid[hi]{aml}{el}{rel}

\olsection{संबंधात्मक प्रतिरूप}

ज्ञानात्मक तर्कशास्त्रों की मूल अर्थवैज्ञानिक धारणा सामान्य मोडल
तर्कशास्त्रों जैसी ही है। संबंधात्मक प्रतिरूप में अब भी जगतों का एक समुच्चय
और ऐसा नियतन होता है जो निर्धारित करता है कि कौन-से !!{propositional variable}
किन जगतों पर ``सत्य'' माने जाते हैं। यदि हम केवल एक कर्ता से संबंधित हैं,
तो सामान्यतः हमारे पास एक ही अभिगम्यता संबंध होता है। किंतु बहु-कर्ता
ज्ञानात्मक तर्कशास्त्र में वह एक संबंध अभिगम्यता संबंधों का समुच्चय बन जाता
है---कर्ता-संकेतों के हमारे समुच्चय~$G$ के प्रत्येक~$a$ के लिए एक संबंध।

किसी \emph{संबंधात्मक प्रतिरूप} में जगतों का एक समुच्चय होता है जो द्विपदी
अभिगम्यता संबंधों---प्रत्येक कर्ता के लिए एक---से जुड़े होते हैं; साथ ही ऐसा
नियतन होता है जो निर्धारित करता है कि कौन-से !!{propositional variable} किन
जगतों पर सत्य हैं।

\begin{defn}
  बहु-कर्ता ज्ञानात्मक भाषा का \emph{प्रतिरूप} एक त्रिक
  $\mModel{M} = \tuple{W, R, V}$ है, जहाँ
  \begin{enumerate}
  \item $W$ ``जगतों'' का अरिक्त समुच्चय है,
  \item प्रत्येक $a \in G$ के लिए ${R}_a$, $W$ पर एक द्विपदी अभिगम्यता
    संबंध है, और
  \item $V$ ऐसा फलन है जो प्रत्येक !!{propositional variable}~$p$ को संभव जगतों
    का समुच्चय $V(p)$ देता है।
  \end{enumerate}
  जब $R_a ww'$ हो, तो हम कहते हैं कि $w'$, $w$ से \emph{$a$ के लिए
    अभिगम्य} है। जब $w \in V(p)$ हो, तो हम कहते हैं कि $p$, $w$ पर
    \emph{सत्य} है।
\end{defn}

इसकी कार्यप्रणाली सामान्य मोडल तर्कशास्त्र जैसी ही है; केवल अधिक अभिगम्यता
संबंध जोड़ दिए गए हैं। किसी दिए हुए कर्ता के अभिगम्यता संबंध की व्याख्या
हम सामान्यतः उसकी सूचना-अवस्था के किसी पक्ष को निरूपित करने वाले संबंध के
रूप में करेंगे। उदाहरण के लिए, हम प्रायः $R_a ww'$ को यह व्यक्त करने वाला
मानते हैं कि $w'$ जगत, $w$ पर कर्ता~$a$ की सूचना के संगत है। दूसरे शब्दों
में, $w$ पर वह कर्ता जगत~$w$ और जगत~$w'$ में अंतर नहीं कर सकता।

\end{document}
